\chapter{Pola wektorowe, przepływy i pochodna Liego}
\label{ch:pola}
\index{pole wektorowe}

W Rozdziale~\ref{ch:rozmaitosci} zbudowaliśmy przestrzenie styczne i wiązkę $TM$.
Teraz wybieramy wektor \emph{w każdym} punkcie, w sposób gładki.
Otrzymujemy pole wektorowe: jednocześnie operator
różniczkowy, kierunek ruchu i infinitezymalny generator lokalnych dyfeomorfizmów.
Z tego jednego pojęcia wynikają nawias Liego, krzywe całkowe, przepływy oraz
pochodna Liego.

Rozdział rozwija rachunek pól oraz geometrię ich przepływów: najpierw
rozwiązujemy równanie krzywej całkowej, a następnie badamy ruch
punktów i sposób przenoszenia funkcji oraz pól wektorowych.
Do dalszej topologii różniczkowej przydadzą się zwłaszcza prostowanie
pola, pola zależne od czasu oraz rozszerzanie izotopii.
Przez $M$ rozumiemy rozmaitość gładką w sensie Definicji~\ref{def:rozmaitosc-gladka},
bez brzegu, chyba że wyraźnie zaznaczono inaczej.
Wszystkie pola i odwzorowania są klasy $C^\infty$.
Uwagi o brzegu i końcach przedziału czasu rozumiemy
według Definicji~\ref{def:rozmaitosc-z-brzegiem-1}
i opisu przestrzeni stycznej w
Sekcji~\ref{sec:rozmaitosci-z-brzegiem-1}.
Tak samo interpretujemy gładkie rodziny na domkniętym przedziale czasu.



\section{Pola wektorowe i derywacje globalne}
\label{sec:derywacje}

W Sekcji~\ref{sec:przestrzen-styczna} utożsamiliśmy wektor $v\in T_pM$ z
\emph{derywacją w punkcie} $v\colon C^\infty(M)\to\R$.
Teraz interesuje nas wybór takiej derywacji gładko w każdym
punkcie $p$. Jest to równoważne jednej derywacji, która każdej
funkcji gładkiej przypisuje \emph{funkcję} gładką.

\begin{definicja}
\emph{Pole wektorowe} $X$ na $M$ to gładki przekrój wiązki stycznej
$TM\to M$, czyli gładkie odwzorowanie $X\colon M\to TM$ spełniające
$X(p)\in T_pM$ dla każdego $p$. Przestrzeń pól oznaczamy $\X(M)=\Gamma(TM)$.
\end{definicja}

W mapie $(U;x^1,\ldots,x^n)$ pole zapisuje się jednoznacznie jako
\begin{equation}\label{eq:lokalne-pole}
 X\big|_U=\sum_{i=1}^{n}X^i\frac{\partial}{\partial x^i},
 \qquad X^i\in C^\infty(U).
\end{equation}
Współczynniki są gładkie; same wektory bazy zmieniają się przy zmianie mapy.
Przykładowo na $\R^2$ pole $X=-y\partial_x+x\partial_y$ jest w każdym punkcie
styczne do okręgu o środku w zerze. W zerze $X_0=0$.

\begin{definicja}
\emph{Derywacją} algebry $C^\infty(M)$ nazywamy $\R$-liniowe odwzorowanie
$D\colon C^\infty(M)\to C^\infty(M)$ spełniające regułę Leibniza
\[
 D(fg)=fD(g)+gD(f),\qquad f,g\in C^\infty(M).
\]
Zbiór takich derywacji oznaczamy $\operatorname{Der}_{\R}(C^\infty(M))$.
\end{definicja}

Pole $X$ działa na funkcje przez $D_X(f)(p):=X_p(f)$.
Współrzędnie $D_X(f)=\sum_iX^i\partial_if$, więc
$D_X$ jest derywacją globalną. Odwrotnie, dla derywacji
globalnej $D$ wartość $v_p(f):=D(f)(p)$ jest liczbą
i spełnia \emph{punktową} regułę Leibniza:
\[
 v_p(fg)=f(p)v_p(g)+g(p)v_p(f).
\]
Zatem $v_p\in T_pM$. Lemat~\ref{lem:lokalnosc} mówi w szczególności,
że $D(f)(p)$ zależy tylko od kiełka $f$ w $p$.
Rozróżniajmy więc $v_p\colon C^\infty(M)\to\R$
oraz $D\colon C^\infty(M)\to C^\infty(M)$.

\begin{twierdzenie}\label{thm:derywacje-pola}
Odwzorowanie
\[
 \X(M)\longrightarrow\operatorname{Der}_{\R}(C^\infty(M)),
 \qquad X\longmapsto(f\longmapsto X(f))
\]
jest izomorfizmem $C^\infty(M)$-modułów.
\end{twierdzenie}
\begin{proof}
Niech $D$ będzie derywacją globalną. Z poprzedniego
rachunku $X_p:=v_p$ jest wektorem stycznym w każdym $p$.
Pozostaje sprawdzić, że $p\mapsto X_p$ jest gładkie.
Wybierzmy wokół dowolnego punktu mapę $(U;x^1,\ldots,x^n)$
i mniejsze otoczenie $V$, którego obraz jest kulą o domknięciu
zawartym w obrazie mapy.
Funkcja odcinająca $\chi$ równa $1$ na $V$
i wsparta w $U$ pozwala przedłużyć każdą współrzędną
do funkcji $\widetilde x^i=\chi x^i\in C^\infty(M)$
(poza $U$ kładziemy ją równą zeru). Dla $p\in V$
lokalność derywacji punktowej daje
\[
 X_p=\sum_i X^i(p)\,\partial_i\big|_p,
 \qquad X^i(p)=v_p(\widetilde x^i)=D(\widetilde x^i)(p).
\]
Każda funkcja $D(\widetilde x^i)$ jest gładka z definicji
derywacji globalnej. Zbiory $V$ pokrywają $M$, zatem
$X$ jest polem gładkim. W drugą stronę gładkie pole $X$
daje derywację $D_X$ przez lokalny wzór
$D_X(f)=\sum_iX^i\partial_if$ i regułę Leibniza.
Obie konstrukcje są odwrotne, a
$D_{fX}=fD_X$ dla $f\in C^\infty(M)$.
\end{proof}

\begin{uwaga}[Dlaczego to jest moduł, a nie zwykle przestrzeń z bazą]
Przestrzeń $\X(M)$ jest modułem nad pierścieniem funkcji: $(fX)_p=f(p)X_p$.
Nie musi mieć jednej globalnej bazy nad $C^\infty(M)$. Na przykład
$TS^2$ nie jest wiązką trywialną; nie da się wybrać dwóch globalnych,
wszędzie liniowo niezależnych pól na sferze. Lokalne bazy
$\partial_{x^1},\ldots,\partial_{x^n}$ zawsze istnieją.
Używamy tu klasycznego twierdzenia o nieistnieniu nigdzie
niezerowego pola na $S^2$; wynika ono również z późniejszego
Twierdzenia~\ref{thm:poincare-hopf-16} i $\chi(S^2)=2$.
\end{uwaga}

\begin{przyklad}[Pole na sferze]
Na $S^2\subset\R^3$ weźmy $X(x,y,z)=(-y,x,0)$.
Iloczyn skalarny $X(x,y,z)\cdot(x,y,z)=0$, zatem $X$ jest styczne
w każdym punkcie. Ma zera na biegunach; poza nimi wskazuje kierunek
obrotu wokół osi $z$. Jest to pole globalne, choć współrzędna kątowa
nie jest globalną mapą na sferze.
\end{przyklad}

\begin{uwaga}[Pole wzdłuż odwzorowania]
Dla gładkiego $F\colon M\to N$ i $X\in\X(M)$ wektory
$\dd F_p(X_p)$ tworzą \emph{pole wzdłuż $F$}.
Precyzyjnie, jest to przekrój wiązki cofniętej
\[
 F^*TN=\{(p,v):p\in M,\ v\in T_{F(p)}N\}\longrightarrow M,
 \qquad(p,v)\longmapsto p.
\]
Mapy na $N$ dają jej lokalne trywializacje nad ich przeciwobrazami
przez $F$; współczynniki przejść są gładkie, bo są złożeniami
przejść $TN$ z $F$. Dla ogólnego $F$ różne punkty $p$ nad jednym
$q=F(p)$ mogą dawać różne wektory, a punkty poza $F(M)$ nie dostają
żadnej wartości. Gdy $F$ jest dyfeomorfizmem, otrzymujemy pole
\[
 (F_*X)_q=\dd F_{F^{-1}(q)}X_{F^{-1}(q)}.
\]
Bardziej ogólnie pole na $N$ może istnieć przy dodatkowej zgodności
wektorów na włóknach $F$, ale nie jest dane automatycznie.
\end{uwaga}

\section{Nawias Liego: mierzenie nieprzemienności}
\label{sec:nawias}

Dwa pola są operatorami na funkcjach. Można je więc składać i wziąć
ich komutator.
\begin{definicja}
\emph{Nawiasem Liego} pól $X,Y\in\X(M)$ nazywamy pole $[X,Y]$ zadane przez
\begin{equation}\label{eq:nawias-def}
 [X,Y](f)=X(Y(f))-Y(X(f)),\qquad f\in C^\infty(M).
\end{equation}
\end{definicja}

Choć $X\circ Y$ jako operator jest na ogół drugiego rzędu, wyrazy z
pochodnymi drugimi znoszą się w różnicy.
\begin{fakt}\label{fakt:nawias-wlasnosci}
$[X,Y]$ jest polem wektorowym. Ponadto nawias jest dwuliniowy nad $\R$,
antysymetryczny i spełnia tożsamość Jacobiego
\[
 [X,[Y,Z]]+[Y,[Z,X]]+[Z,[X,Y]]=0.
\]
Jego reguła Leibniza względem funkcji ma postać
\begin{equation}\label{eq:nawias-leibniz}
 [X,fY]=f[X,Y]+X(f)Y,\qquad [fX,Y]=f[X,Y]-Y(f)X.
\end{equation}
\end{fakt}
\begin{proof}
Pokażmy najpierw, że komutator jest derywacją, choć samo
złożenie dwóch derywacji nią zazwyczaj nie jest.
Dla funkcji $f,g$ dwukrotnie stosujemy regułę Leibniza:
\begin{align*}
[X,Y](fg)
&=X\bigl(Y(f)g+fY(g)\bigr)
  -Y\bigl(X(f)g+fX(g)\bigr)\\
&=X(Yf)g+Y(f)X(g)+X(f)Y(g)+fX(Yg)\\
&\quad-Y(Xf)g-X(f)Y(g)-Y(f)X(g)-fY(Xg)\\
&=[X,Y](f)g+f[X,Y](g).
\end{align*}
Wyrazy mieszane znoszą się dokładnie parami. Komutator
jest liniowy nad $\R$ i spełnia regułę Leibniza,
więc z Twierdzenia~\ref{thm:derywacje-pola} odpowiada
gładkiemu polu wektorowemu. Z definicji natychmiast
$[X,Y]=-[Y,X]$; liniowość w obu argumentach wynika
z liniowości składania operatorów.

Sprawdźmy tożsamość Jacobiego, pokazując wszystkie
wyrazy. Jako operatory na funkcjach mamy
\begin{align*}
[X,[Y,Z]]&=XYZ-XZY-YZX+ZYX,\\
[Y,[Z,X]]&=YZX-YXZ-ZXY+XZY,\\
[Z,[X,Y]]&=ZXY-ZYX-XYZ+YXZ.
\end{align*}
Suma prawych stron jest zerowa: np.\ $XYZ$ znosi się
z $-XYZ$, a analogicznie pozostałe pięć możliwych
kolejności. Ponieważ równość operatorów zachodzi na
każdej funkcji, zachodzi równość pól wektorowych.

Na koniec stosujemy definicję do dowolnej funkcji $g$:
\begin{align*}
[X,fY](g)
 &=X\bigl(fY(g)\bigr)-fY\bigl(X(g)\bigr)
  =X(f)Y(g)+f[X,Y](g),\\
[fX,Y](g)
 &=fX\bigl(Y(g)\bigr)-Y\bigl(fX(g)\bigr)
  =f[X,Y](g)-Y(f)X(g).
\end{align*}
Funkcje rozróżniają pola, zatem są to dokładnie
równości~\eqref{eq:nawias-leibniz}.
\end{proof}

W mapie, jeśli $X=X^i\partial_i$ i $Y=Y^j\partial_j$, to
\begin{equation}\label{eq:nawias-wspolrzedne}
 [X,Y]=\sum_{j=1}^{n}\left(\sum_{i=1}^{n}
 X^i\frac{\partial Y^j}{\partial x^i}-
 Y^i\frac{\partial X^j}{\partial x^i}\right)\partial_j.
\end{equation}
Istotnie, współczynnik przy $\partial_j$ jest równy
$[X,Y](x^j)=X(Y^j)-Y(X^j)$; na funkcjach lokalnych możemy
działać dzięki lokalności derywacji. Tu i dalej $\partial_i$
skraca $\partial/\partial x^i$, a powtórzony indeks w zapisie
takim jak $X^i\partial_i$ oznacza sumowanie po tym indeksie.
Ważny wniosek: $[\partial_i,\partial_j]=0$ w układzie współrzędnych,
ale ogólne pola bazowe, np.\ ramy ruchome, nie muszą komutować.

\begin{przyklad}[Ścinanie płaszczyzny]\label{prz:scinanie}
Niech $X=\partial_x$ i $Y=x\partial_y$ na $\R^2$.
Wtedy $[X,Y]=\partial_y$. Najpierw przesunięcie w kierunku $x$
zmienia prędkość pionową $Y$; wykonanie ruchów w odwrotnej kolejności
daje inny wynik. Pole $Y$ ma zera na osi $x=0$ i nie jest polem
przesunięcia o stałej prędkości.
\end{przyklad}

\begin{przyklad}[Obrót i skalowanie]
Dla $R=-y\partial_x+x\partial_y$ i $E=x\partial_x+y\partial_y$
na $\R^2$ mamy $[R,E]=0$. Obrót i jednolite powiększenie
można wykonywać w dowolnej kolejności. Jeśli zaś
$X=\partial_x$, $Y=x\partial_y$, to wynik jest niezerowy
mimo że w punkcie $(0,0)$ mamy $Y=0$; nawias uwzględnia zmienność
pola w otoczeniu punktu.
\end{przyklad}



\begin{uwaga}[Rozkłady styczne i twierdzenie Frobeniusa]
Rozkładem gładkim stałego rzędu $k$ nazywamy podwiązkę
$\mathcal D\subset TM$, lokalnie rozpiętą przez $k$ gładkich,
punktowo liniowo niezależnych pól. Jest \emph{inwolutywny}, jeśli
$[X,Y]$ jest sekcją $\mathcal D$, gdy $X,Y$ są jej sekcjami.
Twierdzenie Frobeniusa mówi, że dokładnie wtedy lokalnie
istnieją współrzędne $(u^1,\ldots,u^n)$, w których
$\mathcal D=\operatorname{span}(\partial_{u^1},\ldots,\partial_{u^k})$.
Plasterki o stałych $u^{k+1},\ldots,u^n$ są wtedy lokalnymi
podrozmaitościami całkowymi (fragmentami \emph{liści}).
Przykład przeciwny na $\R^3$:
\[
 \mathcal D
 =\operatorname{span}\{\partial_x,\partial_y+x\partial_z\}.
\]
Nawias generatorów jest równy $\partial_z$, który nie należy do $\mathcal D$.
Nie ma zatem powierzchni, których przestrzenie styczne pokrywałyby się
z tym rozkładem. Konieczność warunku jest widoczna już w mapie
plasterkowej: składowe normalne pól stycznych znikają na
podrozmaitości, więc ich pochodne w kierunkach stycznych również
tam znikają; wzór~\eqref{eq:nawias-wspolrzedne} pokazuje,
że nawias jest styczny. W przykładzie $\partial_z$ nie jest
kombinacją generatorów: składowe $x,y$ wymusiłyby zerowe
współczynniki. Dowód dostateczności podamy po twierdzeniu
o prostowaniu pola, w Stwierdzeniu~\ref{prop:frobenius-lokalny}.
\end{uwaga}

\section{Krzywe całkowe i równania różniczkowe na rozmaitości}
\label{sec:krzywe}

\begin{definicja}
\emph{Krzywą całkową} pola $X$ nazywamy krzywą
$\gamma\colon I\to M$ na otwartym przedziale $I\subset\R$ spełniającą
\begin{equation}\label{eq:krzywa-calkowa}
 \gamma'(t)=X_{\gamma(t)}.
\end{equation}
Jeśli $0\in I$ i $\gamma(0)=p$, mówimy o krzywej wychodzącej z $p$.
Jest \emph{maksymalna}, gdy nie daje się przedłużyć jako rozwiązanie
na większy przedział.
\end{definicja}

Współrzędnie, dla $\gamma(t)=(x^1(t),\ldots,x^n(t))$,
warunek~\eqref{eq:krzywa-calkowa} staje się układem równań
różniczkowych zwyczajnych (ODE)
\[
 \frac{\dd x^i}{\dd t}=X^i(x^1(t),\ldots,x^n(t)),\qquad i=1,\ldots,n.
\]
Krzywa nie musi być prostą nawet wtedy, gdy wektory pola rysujemy
jako małe strzałki. W każdym swoim punkcie porusza się w kierunku
aktualnej strzałki; prędkość także jest przez nią określona.

\begin{twierdzenie}[Istnienie, jednoznaczność i zależność od danych]
\label{thm:ode}
Dla każdego $p\in M$ istnieje dokładnie jedna maksymalna krzywa
całkowa $\gamma_p\colon I_p\to M$ pola $X$, gdzie $I_p$ jest otwartym
przedziałem zawierającym $0$ i $\gamma_p(0)=p$.
Rozwiązania zależą gładko od czasu i punktu początkowego tam,
gdzie są określone.
\end{twierdzenie}
\begin{proof}
\emph{Krok 1: konstrukcja w mapie.}
Wybierzmy mapę wokół $p$ i oznaczmy przez $u_0$ współrzędne punktu
$p$. Równanie~\eqref{eq:krzywa-calkowa} przyjmuje postać
$u'(t)=F(u(t))$, gdzie $F$ jest gładkim odwzorowaniem
otwartego zbioru w $\R^n$ do $\R^n$.
Wybierzmy $r>0$ tak, aby domknięta kula
$\overline B_{2r}(u_0)$ leżała w dziedzinie mapy.
Na tej kuli istnieją stałe $A,L\geq 0$ spełniające
$\|F(u)\|\leq A$ oraz $\|\D F(u)\|\leq L$.
Twierdzenie o wartości średniej daje
$\|F(u)-F(v)\|\leq L\|u-v\|$ dla punktów tej kuli.
Wybierzmy $h>0$ tak, aby $hA<r$ i $hL<1$
(gdy $A=0$ lub $L=0$, odpowiedni warunek jest automatyczny).

Dla każdego warunku początkowego $a\in B_r(u_0)$ rozważmy
zupełną przestrzeń krzywych ciągłych
$z\colon[-h,h]\to\overline B_{2r}(u_0)$ z normą supremum.
Operator Picarda
\[
 (\mathcal T_a z)(t)=a+\int_0^t F(z(\tau))\,\dd\tau
\]
odwzorowuje ją w siebie, ponieważ
$\|\mathcal T_a z(t)-u_0\|<r+hA<2r$.
Ponadto
$\|\mathcal T_a z-\mathcal T_a w\|_\infty
 \leq hL\|z-w\|_\infty$.
Z twierdzenia Banacha o punkcie stałym otrzymujemy jedną
krzywą $z_a$ spełniającą równanie całkowe, a po zróżniczkowaniu
$z_a'(t)=F(z_a(t))$. Jest ona gładka w czasie, bo $F$ jest gładkie.

\emph{Krok 2: zależność od punktu początkowego.}
Z równań całkowych dla $z_a$ i $z_b$ oraz z oszacowania
kontrakcji otrzymujemy
\[
 \|z_a-z_b\|_\infty
 \leq \|a-b\|+hL\|z_a-z_b\|_\infty,
 \qquad
 \|z_a-z_b\|_\infty\leq\frac{\|a-b\|}{1-hL}.
\]
To dowodzi najpierw ciągłej zależności, jednak do gładkości
potrzeba więcej. Dla ustalonego $a$ rozwiążmy liniowe
równanie całkowe dla macierzy $J_a(t)$:
\[
 J_a(t)=I+\int_0^t\D F(z_a(\tau))J_a(\tau)\,\dd\tau.
\]
Operator po prawej, działający na ciągłe funkcje o wartościach
macierzowych, również ma stałą kontrakcji $hL<1$.
Istnieje więc jedyne $J_a$; spełnia ono równanie
wariacyjne $J_a'=\D F(z_a)J_a$, $J_a(0)=I$.

Sprawdźmy, że jest to rzeczywiście pochodna względem $a$.
Niech $v$ będzie małe, $\delta=z_{a+v}-z_a$.
Ponieważ $\D^2F$ jest ograniczona na $\overline B_{2r}(u_0)$,
rozwinięcie Taylora na odcinku między $z_a(\tau)$
a $z_{a+v}(\tau)$ daje, z pewną stałą $C$ niezależną
od $v$ i $\tau$,
\[
 \|F(z_a(\tau)+\delta(\tau))-F(z_a(\tau))
       -\D F(z_a(\tau))\delta(\tau)\|
 \leq C\|\delta\|_\infty^2.
\]
Odejmujemy równania całkowe dla $z_{a+v}$, $z_a$
i $J_av$. Dla $R=z_{a+v}-z_a-J_av$ dostajemy
\[
 \|R\|_\infty
 \leq hL\|R\|_\infty+hC\|\delta\|_\infty^2
 \leq hL\|R\|_\infty+
       \frac{hC}{(1-hL)^2}\|v\|^2.
\]
Zatem $\|R\|_\infty=O(\|v\|^2)$ i
$\D_a z_a(t)=J_a(t)$. Dla ciągłości tej pochodnej zauważmy,
że $\|J_a\|_\infty\leq(1-hL)^{-1}$ oraz
\[
 \|J_a-J_b\|_\infty
 \leq\frac{h}{(1-hL)^2}
 \|\D F(z_a)-\D F(z_b)\|_\infty.
\]
Prawa strona dąży do zera dla $b\to a$, dzięki ciągłości
$a\mapsto z_a$ i jednostajnej ciągłości $\D F$ na zwartej kuli.
Dalsze pochodne otrzymujemy w ten sam sposób, ale trzeba
odróżnić ich konstrukcję od formalnego różniczkowania.
Na przykład kandydat na drugą pochodną
$H_a(t)[v,w]$ spełnia
\[
 H_a(t)[v,w]=\int_0^t\bigl(
 \D F(z_a)H_a[v,w]
 +\D^2F(z_a)[J_av,J_aw]\bigr)(\tau)\,\dd\tau.
\]
Najpierw rozwiązujemy to liniowe równanie całkowe dla $H_a$.
Następnie odejmujemy równania dla $J_{a+v}$ i $J_a$ oraz
kandydata $H_a[v,\cdot]$. Wzór Taylora dla $\D F$ i
już dowiedziona różniczkowalność $z_a$ dają dla reszty
oszacowanie $\|R\|_\infty\leq hL\|R\|_\infty+o(\|v\|)$,
a więc $H_a=\D_aJ_a$.
Ogólnie, po $k$ różniczkowaniach złożenia $F(z_a)$
jedynym wyrazem zawierającym nieznaną $k$-tą pochodną $Z_k$
jest $\D F(z_a)Z_k$. Wszystkie pozostałe wyrazy są sumami
iloczynów pochodnych $F$ i już skonstruowanych pochodnych
$Z_j$ z $j<k$. Wyznaczają więc znane ciągłe wymuszenie.
Równanie dla $Z_k$ ma ponownie stałą kontrakcji $hL<1$;
porównanie równań dla $Z_{k-1}(a+v)$ i $Z_{k-1}(a)$
z wzorem Taylora daje tę samą resztę $o(\|v\|)$.
Indukcyjnie dowodzi to istnienia wszystkich pochodnych,
a porównanie równań przy $a,b$ --- ich ciągłości.
Równanie $z_a'=F(z_a)$ zapewnia również ciągłość
pochodnych mieszanych po $t$ i $a$. Otrzymujemy
gładkość $(t,a)\mapsto z_a(t)$.

\emph{Krok 3: przedział maksymalny.}
Dwie krzywe całkowe wychodzące z $p$ są równe w pobliżu zera
na mocy powyższej jednoznaczności. Zbiór czasów, w których
się zgadzają, jest domknięty w części wspólnej ich
przedziałów przez ciągłość obu krzywych. Jest też otwarty:
w punkcie, w którym się zgadzają, stosujemy lokalną
jednoznaczność ponownie. Ponieważ część wspólna jest przedziałem,
krzywe pokrywają się wszędzie, gdzie obie są określone.
Można więc skleić wszystkie rozwiązania przez $p$ na sumie
ich przedziałów. Suma ta jest otwartym przedziałem $I_p$
zawierającym zero; sklejenie daje rozwiązanie, którego żadne
dalsze przedłużenie nie istnieje. Jednoznaczność w mapach
zapewnia, że konstrukcja nie zależy od wyboru współrzędnych.
Gładkość przy innych chwilach niż zero uzyskujemy, zaczynając
lokalny argument od dowolnego punktu tej krzywej; szczegóły
otwartości wspólnej dziedziny podajemy w dowodzie
Twierdzenia~\ref{thm:przeplyw}.
\end{proof}

Jeśli $X_p=0$, to jedyną krzywą wychodzącą z $p$ jest krzywa stała.
Jeśli dwie krzywe całkowe spotkają się, to po odpowiednim przesunięciu
parametru czasu pokrywają się tam, gdzie obie są określone.
Mogą więc tworzyć orbitę okresową, lecz nie mogą przecinać się
transwersalnie jako trajektorie tego samego pola.

\begin{przyklad}[Obrót]
Na $\R^2$ niech $X=-y\partial_x+x\partial_y$. Układ
$x'=-y$, $y'=x$ daje
\begin{equation}\label{eq:obrot}
 \gamma_{(a,b)}(t)=(a\cos t-b\sin t,\ a\sin t+b\cos t).
\end{equation}
Dla $(a,b)\ne(0,0)$ trajektorią jest okrąg o promieniu
$\sqrt{a^2+b^2}$, przebiegany przeciwnie do ruchu wskazówek zegara.
Początek układu jest punktem stałym.
\end{przyklad}

\begin{figure}[!htbp]
\centering
\begin{tikzpicture}[>=Stealth,scale=1.06,line cap=round]
\fill[manifoldblue!4] (-2.5,-2.1) rectangle (2.55,2.1);
\draw[->,manifoldblue!40] (-2.5,0)--(2.55,0) node[right,font=\small] {$x$};
\draw[->,manifoldblue!40] (0,-2.1)--(0,2.1) node[above,font=\small] {$y$};
\draw[manifoldblue!20] (0,0) circle (1.1);
\draw[manifoldblue!32] (0,0) circle (1.78);
\foreach \xx/\yy in {-1.65/-1.1,-1.1/-1.65,0/-1.75,1.25/-1.3,1.72/0,
 1.2/1.25,0/1.65,-1.2/1.2,-1.75/0,-.8/0,.72/.7,-.72/.7,.8/-.7}{
 \draw[->,manifoldgreen!78!black,semithick]
 (\xx,\yy)--++({-.22*\yy},{.22*\xx});}
\draw[manifoldred,very thick,->] (1.78,0)
 arc[start angle=0,end angle=78,radius=1.78];
\fill[manifoldred] (1.78,0) circle (2.4pt)
 node[below right,font=\small] {$p$};
\fill[manifoldred] ({1.78*cos(78)},{1.78*sin(78)}) circle (2.4pt)
 node[above right,font=\small] {$\Phi_t(p)$};
\fill[manifoldblue!80!black] (0,0) circle (2.3pt)
 node[below left,font=\small] {$X_0=0$};
\node[font=\small,manifoldred] at (2.75,1.42) {$\gamma_p$};
\end{tikzpicture}
\caption{Pole obrotowe $X=(-y,x)$: zielone strzałki przedstawiają wartości
pola pomniejszone wspólnym czynnikiem $0.22$,
a czerwona krzywa jest jego krzywą całkową. Strzałka pola
w każdym punkcie jest styczna do trajektorii.}
\label{fig:obrot-krzywa}
\end{figure}

\begin{przyklad}[Zmiana prędkości bez zmiany orbit]
Jeżeli $Y=fX$, gdzie $f>0$, to $X$ i $Y$ wyznaczają te same
\emph{zorientowane ślady} krzywych całkowych, choć czas ich
przebiegania jest inny. Gdy $f$ zmienia znak, orientacja ruchu może się
zmienić; gdy znika, pojawiają się punkty stałe. Dlatego sam obraz
krzywych nie wyznacza pola ani jego przepływu.
Dokładniej, jeśli $\gamma:I\to M$ jest krzywą $X$ z $\gamma(0)=p$,
to funkcja
\[
 s(\tau)=\int_0^\tau\frac{du}{f(\gamma(u))}
\]
ma dodatnią pochodną i jest dyfeomorfizmem $I$ na pewien
przedział $J$. Dla odwrotności $\tau=\tau(s)$ krzywa
$\eta(s)=\gamma(\tau(s))$ spełnia
$\eta'=f(\eta)X_\eta=Y_\eta$. Ta sama konstrukcja z $X=Y/f$
pokazuje równość orbit maksymalnych. Dokładniej, gdyby $\eta$
dała się przedłużyć przez skończony koniec $J$, równanie
$\tau'(s)=f(\eta(s))$ dawałoby tam skończoną wartość $\tau$
i dodatnią pochodną. Odwrotna zmiana czasu przedłużyłaby wtedy
$\gamma$ przez koniec $I$, wbrew jej maksymalności.
Przedziały czasu mogą
jednak być różne: zupełne $X=\partial_x$ po pomnożeniu przez
$f(x)=1+x^2$ daje niezupełne pole z krzywą $\tan s$ przez zero.
\end{przyklad}

\section{Przepływ lokalny, maksymalny i zupełny}
\label{sec:przeplyw}

Krzywe całkowe dla wszystkich punktów początkowych łączą się
w jedno odwzorowanie. W czasie $t$ przesuwamy punkt o tyle, ile
pozwala równanie~\eqref{eq:krzywa-calkowa}.

\begin{definicja}
Niech $\gamma_p\colon I_p\to M$ będzie maksymalną krzywą całkową.
\emph{Dziedziną maksymalnego przepływu} jest zbiór
$\mathcal D_X=\{(t,p)\in\R\times M:t\in I_p\}$, a
\emph{przepływem} pola $X$ odwzorowanie
\[
 \Phi\colon\mathcal D_X\to M,\qquad \Phi(t,p)=\Phi_t(p):=\gamma_p(t).
\]
Pole jest \emph{zupełne}, jeśli $\mathcal D_X=\R\times M$.
\end{definicja}

\begin{twierdzenie}[Podstawowe własności przepływu]
\label{thm:przeplyw}
Zbiór $\mathcal D_X$ jest otwarty, a $\Phi$ jest gładki. Zachodzi
\[
 \Phi_0=\id_M,\qquad
 \Phi_{t+s}(p)=\Phi_t(\Phi_s(p))
\]
dokładnie dla $s\in I_p$ i $t\in I_{\Phi_s(p)}$; przy tym
$I_{\Phi_s(p)}=I_p-s$. Dla $M_t:=\{p\in M:t\in I_p\}$
odwzorowania $\Phi_t\colon M_t\to M_{-t}$ i $\Phi_{-t}$ są wzajemnie
odwrotnymi dyfeomorfizmami między otwartymi zbiorami. Ponadto
\begin{equation}\label{eq:przeplyw-generuje}
 \left.\frac{\dd}{\dd t}\right|_{t=0}\Phi_t(p)=X_p,
 \qquad (\dd\Phi_t)_pX_p=X_{\Phi_t(p)}.
\end{equation}
Jeśli $X$ jest zupełne, $(\Phi_t)_{t\in\R}$ jest jednoparametrową
grupą dyfeomorfizmów.
\end{twierdzenie}
\begin{proof}
\emph{Otwartość i gładkość.}
Krok 1 poprzedniego dowodu daje dla każdego $q\in M$
otoczenie $W_q$ i liczbę $\varepsilon_q>0$, wspólną
dla punktów początkowych z $W_q$, na której lokalne
rozwiązanie istnieje i jest gładkie względem $(u,x)$.
Ustalmy $(t_0,p_0)\in\mathcal D_X$; rozważmy najpierw
$t_0>0$. Krzywa $\gamma_{p_0}$ jest określona
na pewnym otwartym przedziale obejmującym $[0,t_0]$.
Jej obraz $K=\gamma_{p_0}([0,t_0])$ jest zwarty.
Z pokrycia $K$ przez otwarte zbiory $W_q$ wybierzmy
skończone podpokrycie $W_{q_1},\ldots,W_{q_m}$.
Liczba $\varepsilon=\min_i\varepsilon_{q_i}$ jest dodatnia:
z każdego punktu $K$ można rozpocząć lokalne
rozwiązanie przez czas co najmniej $\varepsilon$.
Podzielmy odcinek czasu na
\[
 0=\tau_0<\tau_1<\cdots<\tau_N=t_0
\]
o długościach $\tau_{j+1}-\tau_j<\varepsilon/2$.
Każdy krok wzdłuż $\gamma_{p_0}$ jest wtedy
fragmentem lokalnego rozwiązania z jednego $W_{q_i}$.

Dla $j=0$ lokalny przepływ o czas
$\tau_1-\tau_0$ jest określony i gładki na pewnym
otoczeniu $p_0$. Jego obraz punktu $p_0$ to
$\gamma_{p_0}(\tau_1)$. Ciągłość pozwala zmniejszyć
otoczenie tak, by obrazy pobliskich punktów
należały do dziedziny następnego kroku. Powtarzając
ten argument skończenie wiele razy, otrzymujemy
otoczenie $V$ punktu $p_0$, na którym
$p\mapsto\Phi_{t_0}(p)$ jest gładkim złożeniem
lokalnych rozwiązań. W otoczeniu punktu
$\Phi_{t_0}(p_0)$ można prowadzić rozwiązanie
jeszcze przez wspólny krótki czas $u$.
Po ewentualnym zmniejszeniu $V$ wzór
\[
 \Phi_{t_0+u}(p)
 =\Phi_u\bigl(\Phi_{t_0}(p)\bigr)
\]
wynika z jednoznaczności i definiuje gładkie odwzorowanie
dla $(u,p)$ bliskich $(0,p_0)$. To samo rozumowanie
działa dla $t_0<0$, idąc po krzywej wstecz; przy
$t_0=0$ wystarcza konstrukcja lokalna.
W konsekwencji każde $(t_0,p_0)\in\mathcal D_X$
ma otoczenie zawarte w $\mathcal D_X$, na którym
$\Phi$ jest gładkie. Dowodzi to zarazem otwartości
dziedziny i gładkości przepływu.

\emph{Przesunięcie czasu i prawo grupowe.}
Niech $s\in I_p$ i $q=\Phi_s(p)$. Krzywa
$u\mapsto\gamma_p(s+u)$, określona na $I_p-s$, jest
rozwiązaniem wychodzącym z $q$, więc z jednoznaczności
$I_p-s\subset I_q$. W szczególności $-s\in I_q$.
Krzywa $v\mapsto\gamma_q(v-s)$, określona na $s+I_q$,
jest rozwiązaniem przez $p$ w chwili $v=0$.
Maksymalność $\gamma_p$ daje $s+I_q\subset I_p$.
Oba zawierania dowodzą $I_q=I_p-s$ i
\[
 \Phi_u(q)=\gamma_q(u)=\gamma_p(s+u)=\Phi_{s+u}(p)
 \quad\text{dla każdego }u\in I_q.
\]
Gdy $s\in I_p$, warunki $u\in I_q$ i $s+u\in I_p$
są więc równoważne. Zatem równość składania
obowiązuje dokładnie przy warunkach podanych w tezie.
Szczególny przypadek $s=0$ daje
$\Phi_0(p)=p$.

\emph{Odwrotności i generator.}
Ponieważ $\mathcal D_X$ jest otwarty, każdy $M_t$ jest otwarty.
Jeśli $q=\Phi_t(p)$, to $-t\in I_q$ i prawo składania
daje $\Phi_{-t}(q)=p$. Ten sam argument z zamienionymi
$t$ i $-t$ pokazuje, że $\Phi_t(M_t)=M_{-t}$;
gładkość obu odwzorowań daje dyfeomorfizm.
Pierwszy wzór w~\eqref{eq:przeplyw-generuje} jest równaniem
krzywej całkowej w chwili zero. Dla drugiego różniczkujemy
względem $s$ w zerze równość
$\Phi_t(\Phi_s(p))=\Phi_s(\Phi_t(p))$.
Po lewej stronie reguła łańcucha daje
$(\dd\Phi_t)_pX_p$, po prawej $X_{\Phi_t(p)}$.

Jeżeli $X$ jest zupełne, $M_t=M$ dla każdego $t$.
Wtedy wszystkie $\Phi_t$ są dyfeomorfizmami całej
rozmaitości i
$\Phi_{t+s}=\Phi_t\circ\Phi_s$, $\Phi_0=\id_M$,
$\Phi_t^{-1}=\Phi_{-t}$. Składanie dyfeomorfizmów
jest łączne, a wzór na parametry daje wprost
$\Phi_r\circ(\Phi_s\circ\Phi_t)=\Phi_{r+s+t}
=(\Phi_r\circ\Phi_s)\circ\Phi_t$.
Odwzorowanie $t\mapsto\Phi_t$ jest więc homomorfizmem
$(\R,+)\to\operatorname{Diff}(M)$, a
$(t,p)\mapsto\Phi_t(p)$ jest gładkim działaniem tej
grupy na $M$. Dla pola niezupełnego te same wzory
mają sens tylko w opisanych wyżej otwartych dziedzinach.
\end{proof}

\begin{stwierdzenie}[Grupa przepływu wyznacza swoje pole]
Jeśli gładkie odwzorowanie $H\colon\R\times M\to M$ spełnia
$H_0=\id_M$ i $H_{t+s}=H_t\circ H_s$, to
\[
 X_p:=\left.\frac{\dd}{\dd t}\right|_{t=0}H_t(p)
\]
jest gładkim zupełnym polem wektorowym, a $H_t=\Phi_t^X$
dla wszystkich $t\in\R$.
\end{stwierdzenie}
\begin{proof}
Prawo grupowe daje $H_t\circ H_{-t}=H_{-t}\circ H_t
=H_0=\id_M$. Zatem każdy $H_t$ jest dyfeomorfizmem
o gładkiej odwrotności $H_{-t}$.
Gładkość $H$ zapewnia gładkość przyporządkowania $p\mapsto X_p$.
Z prawa grupowego $H_{t+h}(p)=H_h(H_t(p))$ otrzymujemy,
różniczkując względem $h$ w zerze,
\[
 \frac{\dd}{\dd t}H_t(p)=X_{H_t(p)},\qquad H_0(p)=p.
\]
Zatem $t\mapsto H_t(p)$ jest krzywą całkową $X$ dla każdego
$t\in\R$. Jednoznaczność maksymalnej krzywej daje
$H_t(p)=\Phi_t^X(p)$; istnienie $H_t(p)$ dla wszystkich
czasów dowodzi zupełności pola.
\end{proof}

\begin{figure}[!htbp]
\centering
\begin{tikzpicture}[>=Stealth,scale=1.02,line cap=round]
\fill[manifoldblue!6] (-3.55,-1.42) rectangle (3.65,1.5);
\foreach \c in {-1.05,-.46,.13,.72}{
 \draw[manifoldblue!58,thick]
 plot[domain=-3.2:3.25,samples=55]
 ({\x},{\c+.17*sin(38*\x)});
 \draw[->,manifoldblue!70!black,semithick]
 ({.15},{\c+.17*sin(38*.15)})
 --({.65},{\c+.17*sin(38*.65)});
}
\draw[manifoldred,very thick] (-1.88,{-1.15+.17*sin(38*-1.88)})--(-1.88,{.82+.17*sin(38*-1.88)});
\draw[manifoldgreen!75!black,very thick] (1.68,{-1.15+.17*sin(38*1.68)})--(1.68,{.82+.17*sin(38*1.68)});
\foreach \c in {-1.05,-.46,.13,.72}{
 \fill[manifoldred] (-1.88,{\c+.17*sin(38*-1.88)}) circle (1.9pt);
 \fill[manifoldgreen!75!black] (1.68,{\c+.17*sin(38*1.68)}) circle (1.9pt);
}
\node[font=\small,manifoldred] at (-2.56,1.23) {$S_0$};
\node[font=\small,manifoldgreen!75!black] at (1.7,1.23) {$S_t$};
\node[font=\small,manifoldred] at (-2.18,.06) {$p$};
\node[font=\small,manifoldgreen!75!black] at (2.17,.36) {$\Phi_t(p)$};
\draw[->,very thick,manifoldred] (-.75,1.02)--(.94,1.02)
 node[midway,above,font=\small] {$\Phi_t$};
\end{tikzpicture}
\caption{Przepływ przesuwa punkty wzdłuż krzywych całkowych.
Czerwony przekrój poprzeczny $S_0$ przechodzi w zielony $S_t$;
oba przekroje mają dokładnie te same wartości stałej
$y-A\sin(\kappa x)$, gdzie $A=0.17$ i $\kappa=38\pi/180$.
Są to trajektorie pola $X=(1,A\kappa\cos(\kappa x))$.
Jego przepływ to
$\Phi_s(x,y)=(x+s,y+A[\sin(\kappa(x+s))-\sin(\kappa x)])$;
na rysunku $t=1.68-(-1.88)=3.56$.}
\label{fig:rura-przeplywu}
\end{figure}

\begin{przyklad}[Trzy zupełne przepływy]
Na $\R^n$ pole stałe $X=v$ ma $\Phi_t(x)=x+tv$.
Pole Eulera $E=\sum_i x^i\partial_i$ ma $\Phi_t(x)=e^t x$.
Pole obrotowe z~\eqref{eq:obrot} ma
$\Phi_t(a,b)=(a\cos t-b\sin t,a\sin t+b\cos t)$.
Wszystkie rozwiązania są określone dla całego $t\in\R$.
\end{przyklad}

\begin{przyklad}[Ucieczka w skończonym czasie]\label{prz:niezupelne}
Dla $X=x^2\partial_x$ na $\R$ krzywa przez $x_0$ spełnia
\[
 \Phi_t(x_0)=\frac{x_0}{1-tx_0},\qquad x_0\ne0.
\]
Jeśli $x_0>0$, jej dziedzina to $(-\infty,1/x_0)$; jeśli $x_0<0$,
to $(1/x_0,\infty)$. Dla $x_0=0$ krzywa jest stała i istnieje
dla wszystkich czasów. Pole jest gładkie na całej prostej,
ale nie jest zupełne: nieograniczony wzrost rozwiązania powoduje
ucieczkę do nieskończoności w skończonym czasie.
\end{przyklad}

\begin{twierdzenie}[Kryterium przedłużenia i dwa warunki zupełności]
\label{thm:zupelnosc}
Jeśli maksymalna krzywa całkowa ma skończony prawy koniec $b$, to
opuszcza w pobliżu $b$ każdy zwarty podzbiór $M$.
W konsekwencji każde pole na rozmaitości zwartej jest zupełne.
Każde pole o zwartym nośniku na dowolnej rozmaitości także jest zupełne.
\end{twierdzenie}
\begin{proof}
\emph{Kryterium przedłużenia.}
Niech $I_p=(a,b)$ i $b<\infty$. Załóżmy przeciwnie,
że istnieją zbiór zwarty $K\subset M$ i ciąg
$t_j\nearrow b$ taki, że $\gamma_p(t_j)\in K$.
Po wybraniu podciągu mamy $\gamma_p(t_j)\to q\in K$.
Lokalne twierdzenie o istnieniu w $q$ daje otoczenie
$W\ni q$ i $\varepsilon>0$ takie, że z każdego punktu
$W$ można prowadzić rozwiązanie przez wszystkie czasy
$u\in(-\varepsilon,\varepsilon)$.
Dla dużego $j$ zachodzi $\gamma_p(t_j)\in W$
oraz $b-t_j<\varepsilon/2$.
Krzywa
\[
 u\longmapsto\Phi_u(\gamma_p(t_j)),\qquad
 -\varepsilon<u<\varepsilon,
\]
pokrywa się z $u\mapsto\gamma_p(t_j+u)$ w ich wspólnej
dziedzinie, bo obie zaczynają w $\gamma_p(t_j)$.
Możemy więc skleić je w krzywą określoną także dla
$t> b$, sprzecznie z maksymalnością $I_p$.
Przy skończonym lewym końcu stosujemy ten sam argument
do pola $-X$ i czasu odwróconego.

\emph{Zupełność na zbiorach zwartych.}
Jeśli $M$ jest zwarta, każda trajektoria przez cały czas
leży w jednym zwartym zbiorze $M$. Kryterium wyklucza
oba skończone końce $I_p$, zatem $I_p=\R$.
Teraz niech $K=\supp X$ będzie zwarty.
Jeśli $p\notin K$, to $X_p=0$, więc krzywa
stała w $p$ jest, z jednoznaczności, całą krzywą
maksymalną. Jeśli $p\in K$, krzywa nie może wejść
do $M\setminus K$: od chwili wejścia byłaby stała,
a jednoznaczność, zastosowana także wstecz w czasie,
zmusiłaby ją do bycia stałą już w chwili zero.
W obu przypadkach obraz $\gamma_p$ leży w zwartym
zbiorze $K\cup\{p\}$; kryterium znów daje $I_p=\R$.
\end{proof}

\begin{uwaga}[Nie mylić zupełności z brakiem zer]
Zupełne pole może mieć zera (obrót na płaszczyźnie), a pole bez zer
może być niezupełne, np.\ $X=(1+x^2)\partial_x$ na $\R$.
Dla rozmaitości z brzegiem pole musi dodatkowo być styczne do brzegu,
aby dać dwustronny przepływ zachowujący tę rozmaitość.
Istotnie, dyfeomorfizmy zachowują brzeg, więc różniczkowanie
krzywej $t\mapsto\Phi_t(p)$ dla $p\in\partial M$ daje
$X_p\in T_p\partial M$. Odwrotnie, dla pola stycznego do brzegu
rozwiązujemy równanie również na $\partial M$; jednoznaczność
w lokalnym przedłużeniu pola zapewnia, że jego krzywe przez
brzeg pozostają na brzegu. Krzywe z wnętrza nie mogą go
przekroczyć, bo przy pierwszym spotkaniu jednoznaczność wstecz
wymuszałaby wcześniejszy pobyt na brzegu.
Pole skierowane do wnętrza na brzegu daje zwykle przepływ tylko
dla małych dodatnich czasów, co wykorzystamy do budowania kołnierza.
\end{uwaga}

\section{Geometria przepływu: prostowanie, niezmienniki, komutowanie}
\label{sec:prostowanie}

\begin{definicja}[Orbita i nasycenie przez przepływ]
\index{orbita i nasycenie przez przeplyw@Orbita i nasycenie przez przepływ}
\emph{Orbitą} punktu $p$ pod przepływem pola $X$ jest
$\mathcal O_X(p)=\{\Phi_t(p):t\in I_p\}$.
Dla $A\subseteq M$ jego \emph{nasycenie przez przepływ}
to suma orbit punktów z $A$:
\[
 \operatorname{Sat}_X(A)=\bigcup_{p\in A}\mathcal O_X(p).
\]
Zbiór $A$ jest \emph{nasycony przez $X$}, gdy
$\operatorname{Sat}_X(A)=A$, czyli gdy
$p\in A$ i $t\in I_p$ pociągają za sobą
$\Phi_t(p)\in A$.
\end{definicja}

To jest nasycenie względem relacji ,,leżeć na tej samej
orbicie'' w sensie Definicji~\ref{def:nasycenie}.
Istotnie, gdy $q=\Phi_s(p)$, to
$I_q=I_p-s$ i $\mathcal O_X(q)=\mathcal O_X(p)$
na mocy Twierdzenia~\ref{thm:przeplyw}; orbity rozcinają
więc $M$ na rozłączne klasy. Na przykład dla
$X=\partial_x$ na $\R^2$ nasyceniem
$\{0\}\times(0,1)$ jest cały pas $\R\times(0,1)$.

\begin{twierdzenie}[Twierdzenie o prostowaniu pola]\label{thm:flowbox}
Jeśli $X_p\ne0$, to istnieją współrzędne $(t,z^2,\ldots,z^n)$
w otoczeniu $p$, w których $X=\partial_t$. W tych współrzędnych
lokalny przepływ ma postać
\[
 (t,z^2,\ldots,z^n)\longmapsto(t+s,z^2,\ldots,z^n).
\]
\end{twierdzenie}
\begin{proof}
Wybieramy małą $(n-1)$-wymiarową podrozmaitość $S$
przechodzącą przez $p$ i poprzeczną do $X_p$, oraz parametryzację
$\sigma\colon V\subset\R^{n-1}\to S$ z $\sigma(0)=p$.
Można ją zbudować bez dodatkowego twierdzenia: wybieramy mapę
z $p=0$ i po liniowej zmianie współrzędnych z pierwszą składową
$X_p$ niezerową; bierzemy $S=\{x^1=0\}$ w tej mapie.
Odwzorowanie $F(t,z)=\Phi_t(\sigma(z))$ ma w $(0,0)$
różniczkę, której pierwszą kolumną jest $X_p$, a pozostałe kolumny
rozpinają $T_pS$. Jest więc izomorfizmem. Z twierdzenia o funkcji
odwrotnej $F$ jest lokalnym dyfeomorfizmem. Zmniejszamy jego
dziedzinę do produktu $(-\varepsilon,\varepsilon)\times V_0$.
Ponieważ $\partial_tF(t,z)=X_{F(t,z)}$, we współrzędnych $F^{-1}$
pole ma postać $\partial_t$. Prawo grupowe daje
$\Phi_s(F(t,z))=F(t+s,z)$, dopóki $t$ i $t+s$ leżą
w tym przedziale. To wyjaśnia lokalny zakres wzoru z tezy.
\end{proof}

\begin{uwaga}
Pola nie można wyprostować w zerze $X_p=0$: dyfeomorfizm przenosi
wektor zerowy na wektor zerowy, natomiast $\partial_t$ nigdzie
nie znika. W zerze lokalne trajektorie i zachowanie przepływu
bada się innymi metodami niż prostowanie pola.
\end{uwaga}

\begin{stwierdzenie}[Lokalne twierdzenie Frobeniusa]
\label{prop:frobenius-lokalny}
Gładki rozkład inwolutywny $\mathcal D$ stałego rzędu $k$
ma lokalnie współrzędne, w których jest rozpięty przez pierwsze
$k$ pól współrzędnych.
\end{stwierdzenie}
\begin{proof}
Przeprowadzimy indukcję po $k$, jednocześnie dla wszystkich
wymiarów rozmaitości. Dla $k=0$ teza jest natychmiastowa.
Dla $k\ge1$ wybieramy lokalną bazę $X_1,\ldots,X_k$ rozkładu
i prostujemy niezerowe $X_1$. Mamy więc współrzędne $(t,z)$
w produkcie wokół $(0,0)$, z $X_1=\partial_t$.
Od każdego $X_j$, $j\ge2$, odejmujemy jego składową w kierunku
$\partial_t$. Otrzymane pola $Y_2,\ldots,Y_k$ nie mają tej
składowej, a wraz z $\partial_t$ nadal stanowią bazę $\mathcal D$.
Jeśli $k=1$, dowód jest już zakończony.

Niech $B(t,z)$ będzie macierzą o $k-1$ kolumnach złożonych
ze składowych $Y_2,\ldots,Y_k$ w kierunkach $z$.
Inwolutywność daje $[\partial_t,Y_j]\in\mathcal D$;
ponieważ ten nawias nie ma składowej czasowej, jest kombinacją
samych $Y_2,\ldots,Y_k$. Zatem
\[
 \partial_tB(t,z)=B(t,z)A(t,z)
\]
dla pewnej gładkiej macierzy kwadratowej $A$.
Gładkość współczynników wynika z rozwiązania układu za pomocą
niezerowego minora macierzy $B$, po zmniejszeniu otoczenia.
Rozwiązujemy macierzowe równanie
\[
 \partial_tQ(t,z)=-A(t,z)Q(t,z),\qquad Q(0,z)=I.
\]
Istnienie i gładka zależność od $z$ wynikają z
Twierdzenia~\ref{thm:ode}: dołączamy zmienne z równaniami
$t'=1$, $z'=0$. Po zmniejszeniu produktu $Q$ jest odwracalna,
bo $Q(0,0)=I$. Reguła iloczynu daje
\[
 \partial_t(BQ)=BAQ-BAQ=0,\qquad B(t,z)Q(t,z)=B(0,z).
\]
Kolumny $BQ$ wyznaczają więc nową bazę pól pionowych
niezależnych od $t$, rozpinającą te same kierunki co kolumny $B$.
Na przekroju $t=0$ tworzą rozkład rzędu $k-1$.
Jest on inwolutywny: nawiasy tych pól nie mają składowej
czasowej, należą do $\mathcal D$ i nie zależą od $t$;
reguła Leibniza dla nawiasu rozszerza to na dowolne sekcje.
Z założenia indukcyjnego istnieje zmiana współrzędnych $z$
prostująca ten rozkład. Stosujemy ją niezależnie od $t$.
Ponieważ baza pionowa nie zależy od czasu, w całym produkcie
$\mathcal D$ jest rozpięty przez $\partial_t$ i pierwsze
$k-1$ nowych pól współrzędnych na przekroju. To kończy indukcję.
\end{proof}

\begin{stwierdzenie}[Funkcje i pola niezmiennicze]\label{prop:inwarianty}
Dla zupełnego $X$ funkcja $f$ jest stała na każdej jego trajektorii
wtedy i tylko wtedy, gdy $X(f)=0$. Przepływ zawsze zachowuje
pole, które go generuje: $(\Phi_t)_*X=X$.
Dla innego pola $Y$ następujące warunki są równoważne:
\begin{enumerate}[label=(\roman*)]
\item $[X,Y]=0$;
\item $(\Phi_t)_*Y=Y$ dla wszystkich $t$;
\item lokalne przepływy $X$ i $Y$ komutują tam, gdzie są określone.
\end{enumerate}
W przypadku pól zupełnych ostatni punkt oznacza
$\Phi_t^X\circ\Phi_s^Y=\Phi_s^Y\circ\Phi_t^X$ dla wszystkich $s,t$.
\end{stwierdzenie}
\begin{proof}
Reguła łańcucha daje
$\frac{\dd}{\dd t}f(\Phi_t(p))=X(f)(\Phi_t(p))$.
Druga równość wynika z~\eqref{eq:przeplyw-generuje}.
Dla $Y$ użyjemy wzoru udowodnionego poniżej w
Stwierdzeniu~\ref{prop:lie-ewolucja}:
$\frac{\dd}{\dd t}\Phi_t^*Y=\Phi_t^*[X,Y]$.
Jeśli (i) zachodzi, cofnięte pole ma zerową pochodną
i wartość $Y$ w chwili zero, więc $\Phi_t^*Y=Y$.
Ponieważ $\Phi_t^*Y=(\Phi_{-t})_*Y$, jest to (ii).
Odwrotnie, różniczkowanie (ii) w zerze daje $[X,Y]=0$.
Niech $\Psi_s$ oznacza maksymalny przepływ $Y$. Przy (ii)
dla każdego ustalonego $t$ krzywa $s\mapsto\Phi_t(\Psi_s(p))$
ma pochodną $\dd\Phi_t(Y)=Y$ i zaczyna w $\Phi_t(p)$.
Jednoznaczność daje $\Phi_t(\Psi_s(p))=\Psi_s(\Phi_t(p))$.
Ponieważ $\Phi_t$ jest globalnym dyfeomorfizmem, zastosowanie
tego samego argumentu do $\Phi_{-t}$ pokazuje równość
maksymalnych przedziałów $Y$ przez $p$ i $\Phi_t(p)$.
Wyjaśnia to również zgodność dziedzin w (iii).
Różniczkując (iii) po $s$ w zerze, odzyskujemy
$\dd\Phi_t(Y_p)=Y_{\Phi_t(p)}$, czyli (ii).
\end{proof}

\begin{przyklad}[Pierwsza całka]
Dla pola obrotowego $R=-y\partial_x+x\partial_y$
funkcja $f(x,y)=x^2+y^2$ spełnia $R(f)=0$.
Jej dodatnie poziomice są okręgami, po których biegnie przepływ;
poziomica zerowa to nieruchomy punkt $(0,0)$, a ujemne są puste.
W zagadnieniach dynamiki taka funkcja to \emph{pierwsza całka}.
\end{przyklad}

\section{Pochodna Liego: porównywanie danych w różnych punktach}
\label{sec:liego}

Nie można po prostu odjąć wektorów $Y_{\Phi_t(p)}$ i $Y_p$:
należą do różnych przestrzeni stycznych. Przepływ $X$ dostarcza
\emph{kanonicznego w tej sytuacji} sposobu ich porównania:
różniczka odwrotnego przepływu przenosi pierwszy z nich do $T_pM$.
Otrzymujemy odpowiedź na pytanie, jak obiekt zmienia się
wraz z ruchem generowanym przez $X$.

\begin{definicja}\label{def:pochodna-liego}
Dla funkcji $f$ i pola $Y$ określamy
\begin{align}
 (\Lie_Xf)(p)
 &=\left.\frac{\dd}{\dd t}\right|_{0}f(\Phi_t(p))=X(f)(p),
 \label{eq:lie-f}\\
 (\Lie_XY)_p
 &=\left.\frac{\dd}{\dd t}\right|_{0}
 (\dd\Phi_{-t})_{\Phi_t(p)}Y_{\Phi_t(p)}.
 \label{eq:lie-y}
\end{align}
Definicja jest lokalna: przepływ jest potrzebny tylko dla małego $t$.
Wzór~\eqref{eq:lie-y} zapisujemy krótko jako
$\Lie_XY=\left.\frac{\dd}{\dd t}\right|_0\Phi_t^*Y$.
\end{definicja}


\begin{figure}[!htbp]
\centering
\begin{tikzpicture}[>=Stealth,scale=1.06,line cap=round,line join=round]
% Jedno spójne otoczenie na M; p i q są punktami czerwonej trajektorii.
\filldraw[fill=manifoldblue!8,draw=manifoldblue!70!black,thick]
 (-5.14,-.67) .. controls (-5.15,.52) and (-4.33,1.28) .. (-3.3,1.26)
 .. controls (-1.92,1.22) and (-.88,1.39) .. (.37,1.16)
 .. controls (1.84,1.25) and (4.49,1.43) .. (5.03,.41)
 .. controls (5.52,-.42) and (4.83,-1.3) .. (3.44,-1.43)
 .. controls (1.94,-1.58) and (.42,-1.21) .. (-1.01,-1.44)
 .. controls (-2.84,-1.68) and (-4.84,-1.38) .. cycle;
\node[font=\small,manifoldblue!75!black] at (-4.8,1.51) {$M$};
\coordinate (P) at (-3.65,-.52);
\coordinate (Q) at (3.2,-.48);
% Ciągła krzywa całkowa; wyraźne punkty p, q i kierunek czasu.
\draw[manifoldred!40,thick] (-4.73,-.8)
 .. controls (-4.38,-.71) and (-3.96,-.57) .. (P);
\draw[manifoldred,very thick,postaction={decorate},
 decoration={markings,mark=at position .55 with {\arrow{Stealth}}}]
 (P) .. controls (-2.43,-.1) and (-1.21,-.69) .. (-.1,-.67)
 .. controls (1.13,-.63) and (2.08,-.12) .. (Q);
\draw[manifoldred!40,thick] (Q)
 .. controls (3.81,-.55) and (4.19,-.67) .. (4.67,-.66);
\node[font=\small,manifoldred!85!black] at (.14,-1.08)
 {$\gamma(s)=\Phi_s^X(p)$};
\fill[manifoldred] (P) circle (2.6pt)
 node[below left,font=\small] {$p=\gamma(0)$};
\fill[manifoldred] (Q) circle (2.6pt)
 node[below right,font=\small] {$q=\gamma(t)$};
% Wektory porównywane po sprowadzeniu do T_p M.
\draw[->,manifoldblue,very thick] (P)--(-4.06,.67);
\draw[->,manifoldgold!85!black,very thick] (P)--(-2.78,.78);
\node[font=\small,manifoldblue] at (-4.52,.86) {$Y_p$};
\node[font=\small,manifoldgold!80!black] at (-2.05,.65)
 {$\widetilde Y_t(p)$};
\draw[->,manifoldred,thick] (-4.06,.67)--(-2.78,.78);
\node[font=\small,manifoldred!85!black] at (-3.46,1.01) {$\Delta_t(p)$};
% Wektor Y_q oraz schemat cofnięcia jego wartości, odrębny od trajektorii.
\draw[->,manifoldblue,very thick] (Q)--(3.48,.83);
\node[font=\small,manifoldblue] at (3.98,.89) {$Y_q$};
\draw[->,manifoldgold!85!black,thick,densely dashed]
 (3.48,.83) .. controls (1.93,1.95) and (-.88,2.05) .. (-2.78,.78);
\node[font=\small,manifoldgold!80!black] at (.41,1.92)
 {$(\dd\Phi_{-t}^X)_q$};
\end{tikzpicture}
\caption{Punkty $p=\gamma(0)$ i $q=\gamma(t)=\Phi_t^X(p)$ leżą
na tej samej krzywej całkowej na $M$; czerwona strzałka pokazuje
ruch w czasie od $p$ do $q$. Wektor $Y_q$ cofamy różniczką
przepływu do $T_pM$, otrzymując $\widetilde Y_t(p)$.
Tam porównujemy go z $Y_p$; $\Delta_t(p)$ jest ich różnicą.
Przerywana strzałka przedstawia transport wektora, nie ruch punktu.}
\label{fig:pochodna-liego}
\end{figure}

Dokładnie, oznaczając $q=\Phi_t^X(p)$, mamy
\[
 \widetilde Y_t(p):=(\dd\Phi_{-t}^X)_qY_q\in T_pM,\qquad
 (\Lie_XY)_p=\lim_{t\to0}\frac{\widetilde Y_t(p)-Y_p}{t}.
\]
Dla każdego małego $t$ oba odejmowane wektory leżą w $T_pM$.
Rysunek~\ref{fig:pochodna-liego} pokazuje różnicę $\Delta_t(p)$,
której iloraz przez $t$ ma granicę dla $t\to0$.


\begin{twierdzenie}\label{thm:lie-nawias}
Dla dowolnych pól $X,Y$ zachodzi
\begin{equation}\label{eq:lie-nawias}
 \Lie_XY=[X,Y].
\end{equation}
Zatem pochodna Liego funkcji jest działaniem derywacji,
a pochodna Liego pola jest nawiasem.
\end{twierdzenie}
\begin{proof}
Ustalmy $p\in M$ i $f\in C^\infty(M)$. Pracujemy w otoczeniu $p$
i dla na tyle małych $|t|$, aby lokalny przepływ $\Phi_t=\Phi_t^X$
oraz $\Phi_{-t}$ były określone we wszystkich rozważanych punktach.
Z definicji cofnięcia pola, przy $q=\Phi_t(p)$,
\[
 Z_t(p):=(\Phi_t^*Y)_p=(\dd\Phi_{-t})_qY_q\in T_pM.
\]
Wektory $Z_t(p)$ leżą więc w \emph{jednej} przestrzeni $T_pM$,
a $Z_0(p)=Y_p$. Aby obliczyć ich pochodną, zadziałajmy nimi
na dowolną funkcję $f$. Definicja różniczki daje
\begin{equation}\label{eq:lie-pullback-funkcja}
 Z_t(p)(f)
 =Y_q(f\circ\Phi_{-t})
 =Y(f\circ\Phi_{-t})(\Phi_t(p)).
\end{equation}
Oznaczmy $F(t,r):=f(\Phi_{-t}(r))$ dla $r$ w otoczeniu $p$.
Ponieważ $\frac{\dd}{\dd t}\Phi_{-t}(r)=-X_{\Phi_{-t}(r)}$,
reguła łańcucha daje
\[
 F(0,r)=f(r),\qquad
 \left.\frac{\partial F}{\partial t}\right|_{(0,r)}
 =-X(f)(r).
\]
Niech $H(t,r):=Y_r(F(t,\cdot))$. Wtedy prawa strona
\eqref{eq:lie-pullback-funkcja} jest równa $H(t,\Phi_t(p))$.
Różniczkując to złożenie, musimy uwzględnić zarówno zmianę
funkcji $H$ z czasem, jak i ruch punktu $\Phi_t(p)$:
\begin{align*}
 \left.\frac{\dd}{\dd t}H(t,\Phi_t(p))\right|_{t=0}
 &=\left.\frac{\partial H}{\partial t}\right|_{(0,p)}
   +X_p\bigl(H(0,\cdot)\bigr) \\
 &=-Y_p\bigl(X(f)\bigr)+X_p\bigl(Y(f)\bigr).
\end{align*}
W drugim kroku użyliśmy $H(0,\cdot)=Y(f)$ oraz
$\partial_tH(0,p)=Y_p\bigl(\partial_tF(0,\cdot)\bigr)$.
Zamiana kolejności $Y$ i $\partial_t$ jest dopuszczalna:
w lokalnych współrzędnych $Y=\sum_iY^i(r)\partial_{r^i}$,
współczynniki $Y^i$ nie zależą od $t$, a funkcja $F$ jest gładka.
Wreszcie, ponieważ $Z_t(p)$ jest gładką krzywą w $T_pM$,
jej pochodną można obliczyć po przyłożeniu do $f$. Otrzymujemy
\[
 (\Lie_XY)_p(f)
 =\left.\frac{\dd}{\dd t}Z_t(p)(f)\right|_{t=0}
 =X(Yf)(p)-Y(Xf)(p)=[X,Y]_p(f).
\]
Równość zachodzi dla każdej funkcji gładkiej $f$. Funkcje gładkie
rozróżniają wektory styczne: wystarczy wziąć funkcje, które
w pobliżu $p$ pokrywają się ze współrzędnymi lokalnymi.
Zatem $(\Lie_XY)_p=[X,Y]_p$; punkt $p$ był dowolny.
\end{proof}

\begin{stwierdzenie}\label{prop:lie-ewolucja}
Na wspólnych dziedzinach lokalnego przepływu pola $X$ zachodzi
\[
 \frac{\dd}{\dd t}\Phi_t^*Y=\Phi_t^*[X,Y].
\]
\end{stwierdzenie}
\begin{proof}
Dla dyfeomorfizmu $F$ cofnięcie pola oznacza $F^*Y=(F^{-1})_*Y$.
Reguła łańcucha daje $(F\circ G)^*=G^*F^*$.
Prawo grupowe przepływu pozwala zatem napisać
$\Phi_{t+h}^*Y=\Phi_t^*(\Phi_h^*Y)$.
Ustalamy punkt i dostatecznie małe wspólne otoczenie, aby
wszystkie te odwzorowania były określone dla $h$ bliskiego zeru.
Różniczkujemy po $h$ w zerze: liniowa mapa cofająca przez
ustalone $\Phi_t$ nie zależy od $h$, a
$\left.\partial_h\right|_0\Phi_h^*Y=[X,Y]$
na mocy Twierdzenia~\ref{thm:lie-nawias}.
\end{proof}

\begin{uwaga}[Ten sam rachunek we współrzędnych]
W lokalnych współrzędnych, zapisując pola jako kolumny funkcji,
gładkość przepływu daje przy $t\to0$
\[
 \Phi_t(x)=x+tX(x)+o(t),\qquad
 \dd\Phi_t|_x=I+t\,\dd X|_x+o(t),
\]
a stąd $\dd\Phi_{-t}|_{\Phi_t(x)}=(\dd\Phi_t|_x)^{-1}
=I-t\,\dd X|_x+o(t)$ oraz
$Y(\Phi_t(x))=Y(x)+t\,\dd Y|_xX(x)+o(t)$.
Po pomnożeniu obu wyrażeń otrzymujemy
\[
 (\Phi_t^*Y)(x)=Y(x)+t\bigl(\dd Y|_xX(x)-\dd X|_xY(x)\bigr)+o(t).
\]
Zatem składowa $j$ pochodnej wynosi
$\sum_i(X^i\partial_iY^j-Y^i\partial_iX^j)$,
co zgadza się ze wzorem na nawias Liego.
\end{uwaga}

\begin{przyklad}[Obliczenie w polu przesunięć]
Dla $X=\partial_x$ i $Y=x\partial_y$ mamy
$\Phi_t(x,y)=(x+t,y)$. W punkcie $q=(x+t,y)$ jest
$Y_q=(x+t)\partial_y$, a cofająca różniczka przesunięcia
działa identycznościowo na wektory. Zatem
\[
 ((\Phi_t)^*Y)_{(x,y)}=(x+t)\partial_y,
 \qquad \Lie_XY=\partial_y.
\]
Nie chodzi o różnicę położeń punktów, tylko o zmianę pola
po przeniesieniu obu wektorów do tej samej przestrzeni stycznej.
\end{przyklad}

\begin{figure}[!htbp]
\centering
\begin{tikzpicture}[>=Stealth,scale=1.12,line cap=round,line join=round]
% Cztery kroki komutatora dla X=\partial_x, Y=(1+x)\partial_y.
\coordinate (P) at (-2,-.68);
\coordinate (A) at (.7,-.68);
\coordinate (B) at (.7,1.04);
\coordinate (C) at (-2,1.04);
\coordinate (Q) at (-2,.04);
\fill[manifoldblue!5] (-2.65,-1.15) rectangle (1.38,1.52);
\draw[manifoldblue!22,thin] (-2.65,-1.15) rectangle (1.38,1.52);
\draw[->,very thick,manifoldblue] (P)--(A)
 node[midway,below,font=\small] {$\Phi_t^X$};
\draw[->,very thick,manifoldgreen] (A)--(B)
 node[midway,right,font=\small] {$\Phi_t^Y$};
\draw[->,very thick,manifoldblue] (B)--(C)
 node[midway,above,font=\small] {$\Phi_{-t}^X$};
\draw[->,very thick,manifoldgreen] (C)--(Q)
 node[midway,left,font=\small] {$\Phi_{-t}^Y$};
\fill[manifoldred] (P) circle (2.5pt) node[below left,font=\small] {$p$};
\fill[manifoldred] (Q) circle (2.5pt) node[left,font=\small] {$C_t(p)$};
\draw[->,thick,manifoldred,densely dashed] (-1.66,-.68)--(-1.66,.04)
 node[midway,right,font=\small] {$t^2[X,Y]_p$};
\node[font=\footnotesize,align=center] at (3.7,.52)
 {$C_t=\Phi_{-t}^Y\circ\Phi_{-t}^X$\\[2pt]
 $\circ\Phi_t^Y\circ\Phi_t^X$};
\node[font=\footnotesize,align=center,text width=3.4cm] at (3.75,-.6)
 {Końcowy punkt leży o rząd $t^2$ ponad początkiem.};
\end{tikzpicture}
\caption{Komutator małych przepływów. Po przejściu kolejno wzdłuż
$X,Y,-X,-Y$ droga zwykle się nie zamyka; wyraz przemieszczenia rzędu $t^2$ jest wyznaczony przez nawias
$[X,Y]$. Jeśli nawias znika, pierwszy niezerowy wyraz może mieć
wyższy rząd. Rysunek schematycznie ilustruje przykład
$X=\partial_x$, $Y=(1+x)\partial_y$.}
\label{fig:komutator-przeplywow}
\end{figure}

Jeśli $C_t=\Phi_{-t}^Y\circ\Phi_{-t}^X\circ\Phi_t^Y\circ\Phi_t^X$, to
w dowolnej mapie wokół $p$ (i dla małego $t$)
\[
 x(C_t(p))=x(p)+t^2\,\dd x_p([X,Y]_p)+O(t^3).
\]
W przykładzie $X=\partial_x$, $Y=(1+x)\partial_y$ i $p=(0,0)$
można to policzyć dokładnie: $p\to(t,0)\to(t,t+t^2)
\to(0,t+t^2)\to(0,t^2)$. Geometria rysunku zgadza się więc ze
znakiem konwencji~\eqref{eq:nawias-def}.

Uzasadnijmy także wzór ogólny, zamiast wnioskować z jednego przykładu.
W ustalonej mapie przepływ gładkiego pola $Z$ ma rozwinięcie
\[
 \Phi_t^Z(x)=x+tZ(x)+\tfrac12t^2 DZ(x)Z(x)+O(t^3),
\]
bo druga pochodna krzywej całkowej w zerze to $DZ(x)Z(x)$.
Oznaczmy kolejne punkty czterech kroków przez $x_1,\ldots,x_4$.
Rozwijając pola w punktach pośrednich, dostajemy
\begin{align*}
 x_1&=x+tX+\tfrac12t^2 DX\,X+O(t^3),\\
 x_2&=x+t(X+Y)+t^2(\tfrac12 DX\,X+DY\,X+\tfrac12 DY\,Y)+O(t^3),\\
 x_3&=x+tY+t^2(DY\,X+\tfrac12 DY\,Y-DX\,Y)+O(t^3),\\
 x_4&=x+t^2(DY\,X-DX\,Y)+O(t^3).
\end{align*}
Wszystkie pola i ich macierze pochodnych po prawej stronie
są obliczone w $x$. Równość $DY\,X-DX\,Y=[X,Y](x)$
wynika ze wzoru~\eqref{eq:nawias-wspolrzedne}.
Reszty $O(t^3)$ są rozumiane w tej mapie, dla punktów w
dostatecznie małym otoczeniu i małego czasu, gdy złożenia istnieją.


\begin{uwaga}[Znaczenie znaku]
W~\eqref{eq:lie-y} wektor cofamy przez $\Phi_{-t}$.
Przy tej konwencji $\Lie_XY=[X,Y]$. Gdyby transportować
$Y_p$ w przód i porównywać z $Y_{\Phi_t(p)}$ w punkcie
$\Phi_t(p)$ w odwrotnej kolejności, otrzymalibyśmy zapis
z przeciwnym znakiem. Przed rachunkiem warto zawsze ustalić,
które dwa wektory są porównywane i w jakim punkcie.
\end{uwaga}

\section{Pola zależne od czasu i izotopie}
\label{sec:izotopie}

Przepływ pola stałego w czasie spełnia lokalne prawo grupowe;
dla pola zupełnego tworzy jednoparametrową grupę dyfeomorfizmów.
W topologii różniczkowej często trzeba deformować obiekty według
zaplanowanego ruchu, którego prędkość może zmieniać się z czasem.

\begin{definicja}
\emph{Pole zależne od czasu} to rodzina $X_t\in\X(M)$,
$t\in[0,1]$, gładka łącznie w $(t,p)$ (na końcach przedziału
w sensie lokalnego gładkiego przedłużenia). Jej krzywe całkowe spełniają nieautonomiczne równanie
$\gamma'(t)=(X_t)_{\gamma(t)}$. Odwzorowanie ewolucji
$F_{t,s}(p)$ przenosi położenie w chwili $s$ do położenia
w chwili $t$, tam gdzie rozwiązanie istnieje.
\end{definicja}

Dodanie czasu jako zmiennej sprowadza to równanie do autonomicznego
pola $\widehat X=\partial_t+X_t$ na czasoprzestrzeni.
Krzywa tego pola startująca w $(s,p)$ ma postać
$(s+u,F_{s+u,s}(p))$. Twierdzenie~\ref{thm:ode}, zastosowane
lokalnie także do przedłużeń przy końcach przedziału, daje
gładkość ewolucji łącznie w $(t,s,p)$. Jednoznaczność daje prawo
\begin{equation}\label{eq:ewolucja}
 F_{s,s}=\id,\qquad F_{u,t}\circ F_{t,s}=F_{u,s},
 \qquad F_{s,t}=F_{t,s}^{-1}.
\end{equation}
Równości obowiązują tam, gdzie odpowiednie rozwiązania istnieją;
$F_{t,s}$ jest dyfeomorfizmem swojej otwartej dziedziny na obraz,
z odwrotnością $F_{s,t}$.
Zwykle $F_{t,0}$ \emph{nie} spełnia
$F_{t+s,0}=F_{t,0}\circ F_{s,0}$, bo samo pole się zmienia.
Jeżeli nośniki wszystkich $X_t$ mieszczą się w jednym zwartym
$K\subset M$, ewolucja istnieje dla całego przedziału $[0,1]$:
poza $K$ krzywe stałe są rozwiązaniami, a jednoznaczność
w obu kierunkach czasu nie pozwala trajektorii z $K$ wyjść
poza $K$. Trajektoria startująca w $p$ pozostaje więc w
$K\cup\{p\}$. Jej podniesienie do czasoprzestrzeni leży
w zwartym $[0,1]\times(K\cup\{p\})$. Argument przedłużania
z dowodu Twierdzenia~\ref{thm:zupelnosc}, z lokalnymi przedłużeniami
pola na końcach przedziału, wyklucza skończony czas zakończenia
rozwiązania przed dojściem do wymaganej chwili. Istotny jest
\emph{wspólny} zwarty nośnik, a nie tylko zwartość każdego z nośników osobno.

\begin{definicja}
Gładka rodzina dyfeomorfizmów $h_t\colon M\to M$ z
$h_0=\id_M$ to \emph{izotopia otoczenia}.
Gładka rodzina osadzeń $i_t\colon N\hookrightarrow M$ to
\emph{izotopia osadzeń}. Drugie pojęcie dotyczy tylko ruchu
podrozmaitości, pierwsze przesuwa całe otoczenie.
Osadzenie oznacza tu zanurzenie w sensie Rozdziału~\ref{ch:rozmaitosci},
czyli immersję będącą homeomorfizmem na obraz.
Nośnikiem izotopii otoczenia nazywamy domknięcie zbioru
$\{p\in M:h_t(p)\ne p\text{ dla pewnego }t\in[0,1]\}$.
\end{definicja}

Każda izotopia otoczenia jest generowana przez pole zależne od
czasu: dla $q=h_t(p)$ kładziemy
\begin{equation}\label{eq:generator-izotopii}
 (X_t)_q=\left.\frac{\partial h_t(p)}{\partial t}
 \right|_{p=h_t^{-1}(q)}.
\end{equation}
Gładkość tego pola wymaga także gładkości $(t,q)\mapsto h_t^{-1}(q)$.
Wynika ona z twierdzenia o funkcji odwrotnej zastosowanego do
$H(t,p)=(t,h_t(p))$: różniczka jest blokowo trójkątna z
odwracalnymi blokami $1$ i $\dd h_t$. Przy $t=0,1$ stosujemy
lokalne przedłużenia; pierwsza współrzędna pozostaje czasem.
Jednoznaczność rozwiązań daje $h_t=F_{t,0}$.
W szczególności pole zależne od czasu
o wspólnym zwartym nośniku buduje izotopię otoczenia
stałą poza tym nośnikiem.

\begin{twierdzenie}[Rozszerzenie izotopii, wersja zwarta]
\label{thm:rozszerzenie-izotopii}
Niech $M$ będzie rozmaitością bez brzegu,
$N$ jej zwartą podrozmaitością bez brzegu,
a $i_t\colon N\hookrightarrow M$, $t\in[0,1]$, gładką
izotopią osadzeń. Istnieje izotopia otoczenia $h_t$ taka, że
$h_t\circ i_0=i_t$. Jej nośnik można umieścić w dowolnym
otwartym otoczeniu zwartego śladu
$\bigcup_{t\in[0,1]}i_t(N)$.
\end{twierdzenie}
\begin{proof}
\emph{Krok 1: pole na śladzie izotopii.}
Zbudujmy \emph{ślad w czasoprzestrzeni}
\[
 \mathcal T:=\{(t,i_t(x)):t\in[0,1],\ x\in N\}
 \subset[0,1]\times M.
\]
Odwzorowanie $I(t,x)=(t,i_t(x))$ jest różnowartościowe:
z równości pierwszych współrzędnych mamy ten sam $t$,
a różnowartościowość $i_t$ daje ten sam $x$.
Jego różniczka jest injektywna: jeśli
$\dd I_{(t,x)}(a,v)=0$, to pierwsza współrzędna daje
$a=0$, a następnie $\dd (i_t)_x(v)=0$ wymusza $v=0$.
Ponieważ $[0,1]\times N$ jest zwarta, $I$ jest osadzeniem
(ciągła bijekcja z przestrzeni zwartej na podprzestrzeń
Hausdorffa ma ciągłą odwrotność). Zatem $\mathcal T$
jest zwartą podrozmaitością z brzegiem w $[0,1]\times M$.
Wyjaśnijmy postać lokalną także dla $t=0,1$. W lokalnych
współrzędnych $u$ na $N$ wybieramy $k=\dim N$ składowych $i_t(u)$,
których różniczka względem $u$ jest odwracalna. Odwzorowanie
$(t,u)\mapsto(t,i_t^1(u),\ldots,i_t^k(u))$ jest lokalnym
dyfeomorfizmem. Pozostałe składowe śladu są więc wykresem
$w=b(t,z)$. Zmiana $(t,z,w)\mapsto(t,z,w-b(t,z))$ prostuje
ślad i zachowuje czas, a zatem również brzeg przy $t=0,1$.
Homeomorficzność $I$ na obraz pozwala zmniejszyć otoczenie tak,
aby nie zawierało innych fragmentów śladu.
Wzdłuż niej mamy dobrze określony \emph{pionowy} wektor
\[
 V(t,i_t(x)):=\frac{\partial i_t(x)}{\partial t}
 \in T_{i_t(x)}M.
\]
Jest gładki, bo $I$ ma gładką odwrotność na $\mathcal T$.

\emph{Krok 2: przedłużenie o małym nośniku.}
Niech $O\subset M$ będzie wskazanym otwartym otoczeniem
rzutu $\mathcal T$ na $M$. Wokół każdego punktu
$\mathcal T$ wybieramy lokalne współrzędne na
$[0,1]\times M$, w których $\mathcal T$ ma postać
normalną opisaną powyżej; otoczenia wybieramy wewnątrz
$[0,1]\times O$. W lokalnej trywializacji $TM$
składowe wektora $V$ są gładkimi funkcjami współrzędnych
na $\mathcal T$; przedłużamy je, nadając im te same
wartości w pobliskich kierunkach normalnych.
Otrzymujemy lokalne pola pionowe $V^{(j)}$, których
ograniczenia do $\mathcal T$ są równe $V$.
Zwartość $\mathcal T$ pozwala wybrać skończenie wiele
takich otoczeń i nieujemne gładkie funkcje odcinające $\chi_j$
o zwartych nośnikach w tych otoczeniach, z
$\sum_j\chi_j>0$ na $\mathcal T$.
Na otwartym zbiorze $U=\{\sum_j\chi_j>0\}$ kładziemy
$Y=(\sum_j\chi_j V^{(j)})/(\sum_k\chi_k)$,
przedłużając każde $\chi_j V^{(j)}$ przez zero poza jego otoczenie.
Na $\mathcal T$ wszystkie składniki $V^{(j)}$ są równe
$V$, więc $Y|_{\mathcal T}=V$.
Wybieramy jeszcze funkcję odcinającą $\rho$ równą $1$ w
otoczeniu zwartej $\mathcal T$, o zwartym nośniku
w $U\subset[0,1]\times O$, i mnożymy przez nią $Y$.
Taka funkcja powstaje ze skończonego pokrycia $\mathcal T$
otoczeniami z funkcjami $0\le\rho_j\le1$, równymi $1$
na mniejszych otoczeniach: $\rho=1-\prod_j(1-\rho_j)$.
Przy końcach czasu wystarczą ograniczenia zwykłych funkcji
odcinających do półprzestrzeni. Warunek $\supp\rho\subset U$
zapewnia, że nie dzielimy przez zero.
Przedłużenie przez zero definiuje gładkie pole $X_t$
na całym $M$, gładkie w $t$, którego wszystkie nośniki
leżą w jednym zwartym $K\subset O$ (rzucie $\supp\rho$ na $M$)
i które spełnia
$(X_t)_{i_t(x)}=\partial_t i_t(x)$.

\emph{Krok 3: całkowanie.} Każda krzywa
$t\mapsto i_t(x)$ spełnia równanie
$\dot q(t)=(X_t)_{q(t)}$ z warunkiem początkowym
$q(0)=i_0(x)$. Lokalna jednoznaczność dla pola zależnego
od czasu wynika z tego samego argumentu co
Twierdzenie~\ref{thm:ode}, po dodaniu czasu jako zmiennej.
Wspólny zwarty nośnik $K$ wyklucza ucieczkę rozwiązania
w skończonym czasie, więc ewolucja $F_{t,s}$ istnieje
dla $s,t\in[0,1]$. Z jednoznaczności
$F_{t,0}(i_0(x))=i_t(x)$. Kładąc $h_t:=F_{t,0}$,
otrzymujemy izotopię dyfeomorfizmów:
jej odwrotnością w chwili $t$ jest $F_{0,t}$.
Poza $K$ pole $X_t$ znika, zatem punkty tam pozostają
nieruchome i nośnik $h_t$ leży w $O$.
\end{proof}

\begin{przyklad}[Przesuwanie punktu]
Gdy $i_t\colon\{*\}\hookrightarrow M$ opisuje gładką drogę
$p_t$ w spójnej składowej $M$, rozszerzenie izotopii daje
rodzinę dyfeomorfizmów $h_t$ z $h_t(p_0)=p_t$.
Można ją zbudować konkretnie, przedłużając prędkość
$\dot p_t$ do pola wspieranego w małym otoczeniu całej drogi.
To narzędzie pozwala przesuwać lokalne konstrukcje
bez zmieniania odległych części rozmaitości.
\end{przyklad}

\begin{uwaga}[Kołnierz brzegu]
Jeśli $M$ ma zwarty brzeg, wybieramy pole $X$, które na
$\partial M$ jest skierowane do wnętrza i poprzeczne do brzegu.
Jego mały dodatni przepływ daje dyfeomorfizm na obraz
$\partial M\times[0,\varepsilon)\to M$,
$(p,t)\mapsto\Phi_t(p)$, zwany \emph{kołnierzem}.
Pole takie otrzymujemy z lokalnych pól skierowanych do wnętrza
przez nieujemne funkcje odcinające; dodatnia składowa normalna
pozostaje dodatnia przy ich sumowaniu. W mapach półprzestrzeni
przedłużamy pole za brzeg, aby zastosować lokalne twierdzenie
o przepływie. Jego dodatnie trajektorie pozostają początkowo w $M$.
Różniczka $(v,a)\mapsto v+aX_p$ jest izomorfizmem w $(p,0)$.
Twierdzenie o funkcji odwrotnej, z uwzględnieniem znaku składowej
normalnej, daje lokalny dyfeomorfizm półotoczeń. Zwartość brzegu
pozwala wybrać wspólny czas istnienia i lokalnej odwracalności.
Pozostaje globalna injektywność dla małego czasu. Gdyby żadna
wartość $\varepsilon$ jej nie dawała, istniałyby różne pary
$(p_j,t_j),(q_j,s_j)$ o tym samym obrazie i $0\le t_j,s_j<1/j$.
Po przejściu do podciągu $p_j\to p$, $q_j\to q$; ciągłość
daje $p=q$. Obie pary leżą wtedy w jednym otoczeniu lokalnej
injektywności przy $(p,0)$, sprzeczność. Obraz jest otwartym
otoczeniem brzegu w $M$, bo mapa jest lokalnym dyfeomorfizmem
półotoczeń. Kołnierz pozwala opisywać otoczenie brzegu
współrzędnymi niezależnymi od punktu na brzegu.
\end{uwaga}

\section*{Dalsza lektura}
\addcontentsline{toc}{section}{Dalsza lektura}

\begin{itemize}[leftmargin=*]
\item \href{https://doi.org/10.1007/978-1-4419-9982-5}
      {J. M. Lee, \emph{Introduction to Smooth Manifolds}, wyd. 2},
      rozdziały 8--9 i 12: pola, przepływy, pochodna Liego
      i rozszerzanie izotopii.
\item \href{https://people.dm.unipi.it/benedett/MILNOR-TDVPOINT.pdf}
      {J. Milnor, \emph{Topology from the Differentiable Viewpoint}},
      §4 i §6: gładka izotopia oraz pola wektorowe.
\end{itemize}
\clearpage
